\chapter{Point Estimations}
We have thus far talked about functional inferences, where we try to find a function that expresses some degree of "plausibility" for different values of $\theta$ starting from the observations of $x$ and its distribution $p(x;\theta)$. 

We now pivot to something entirely different, and our aim will be, starting from the measurements of $x$\sidenote{That is all we can do, after all, measure x! Remember, our job as data analysts is to infer something from the measured data.}, to extract a single value for the parameter $\theta$, attributing to that value some uncertainty. We have seen this many times in our lives, and usually write it in the synthetic form:
	\[
	\theta = \hat{\theta} \pm e
	\]
Where $\hat{\theta}$ is the \emph{point estimation} of the parameter $\theta$. The idea is that this point estimation gives us the "best estimate" for the parameter $\theta$, whatever that means. 

Remember, we can't measure parameters, so usually when we read that a parameter (let's say, the lifetime of the muon) was \emph{measured}, it really means that it was \emph{estimated through a point estimation} and the value of that estimation is provided, usually with an associated statistic uncertainty. 

Our job during this chapter will be to explore the different kinds of point estimators and to understand their properties so that we can decide, in the course of our data analyst careers, which tool is right for which job. Unfortunately, there is no "best way" that fits all of the situations we might find ourselves in, and the expert analyst is the one that knows how to justify their choice in a convincing matter.\sidenote{Note: not the one who makes the \emph{best} choice, as often it's unclear which properties we can afford to abandon in favor of others.}

\section{Bayesian point estimator}
Let's, as we have done before, discuss the Bayesian case first, where our parameter $\theta$ follows a pdf. In this case, it's natural to start from the posterior as described in (\ref{def:posterior}), $p(\theta|x_0)$, and take its maximum as the estimator $\hat{\theta}$ for $\theta$. 

This is a natural yet unsatisfying approach - I know everything I could ever hope to know about $\theta$ through the full knowledge of the posterior and I waste the whole shape of the distribution only considering the maximum? We can definitely do better than that, and often when we are fortunate enough to find ourselves in this situation we can choose to abandon point estimates entirely in favor of other techniques that fully use the information contained in the posterior. 

For example, usually we can talk about the problem in terms of minimizing some sort of risk function $r(\tilde{\theta},\theta)$, which for example can have the form $r = (\tilde{\theta}-\theta)^2$ (quadratic risk). We intend that
	\[
	\tilde{\theta} = \E_\theta[\theta] = \int\theta\cdot p(\theta|x_0)\dd \theta
	\]

From this we can see that taking the maximum of the posterior is equivalent to using the risk function $r(\tilde{\theta},\theta)=P(\tilde{\theta}\neq\theta)$, which makes sense to use if $\theta$ is discrete. 
\todo{example of quadratic and maximum risk?}

To be honest, this formalism will not be looked further upon, and is rarer to find with respect to the "classic" frequentist approach, which is the one that is blindly intended when not specified, and on which we will spend the rest of the chapter.
\todo{Bayesian point estimation is used in Compito 16/07/2018 Problema 1 Punto 6. We might want to insert that extract as an example since we never displayed (even in the lecture) how to perform such point estimation.}

\section{The classic point estimator}

What is an "estimate of $\theta$"? Well, really any statistic $s(x)$ can be an estimate of $\theta$, we will baptize it \emph{estimator} and denote it with the symbol $\hat{\theta}$. What we will try to answer in this section is how \emph{good} a certain estimator $\hat{\theta}$ is.

For example, I can choose $\hat{\theta}=42$\sidenote{It is the answer to the life, the universe, and everything after all!}. We can imagine that this estimator will not be particularly good nor interesting, and it's not unreasonable to request that the estimator for $\theta$ have something to do with $\theta$ itself.

Let's take a look at the situation we find ourselves in. Differently from the Bayesian case, we want to estimate a parameter that is not governed by a probability function, much less defined on a metric, and try to build a formalism that quantitatively tells us how "close" we are to estimating the parameter itself.\sidenote{Naturally, we can't talk about the "probability of guessing" $\theta$, as the phrase itself is not well defined, and we can't attribute a "real" value to the parameter any better than we can estimate it.} So how do we approach this daunting task?

The answer lies in the fact that, much like the measurements, the estimate can be repeated an infinite number of times. Therefore, there exists the pdf of $s(x)$ (and the estimator of $N$ measurements can be generally written as $s_N(\{x_i\})$). 
\begin{definition}[Consistent estimator]
	The estimator $s_N(\{x_i\})$ is a \emph{consistent} estimator of $\theta$ if
	\[
	\forall \varepsilon > 0: \lim_{N\to\infty} P(|s_N(x)-\theta|>\varepsilon)=0
	\]
	which, we hope you remember, means that $s_N$ \emph{converges in probability to} $\theta$.
\label{def:consistent_estimator}
\end{definition}

Generally, we believe it's more clear to denote the estimator for $\theta$ as $\hat{\theta}$, and so from now on we'll use this notation, remembering that $\hat{\theta}$ is a statistic (and so is a function of the data) that admits pdf.

This property is so important that if $\hat{\theta}$ is not at least consistent for $\theta$ then  it is not even considered an estimator to begin with; frankly, it's non-negotiable that an estimator be at least consistent. For this reason, if we ever talk about any sort of estimator, this property is silently assumed. 

With this property we can re-state the LLN \ref{subsec:LLN} as "the arithmetic mean is a consistent estimator of the mean $\E[x]$ of the distribution." We have already stated this, but it's helpful to remember if things get confusing; when we need a concrete example of an estimator let's think of the sample mean as a (consistent) estimator for the mean of a distribution.

At this point, it becomes interesting to consider the variance of the estimator\sidenote{Remember, $\hat{\theta}$ admits a pdf, so talking about the variance makes sense}, and since we know that it converges to the parameter at high statistics, we want to quantify "how precisely" it does so. In fact, when we write 
	\[
	\theta = \hat{\theta} \pm e
	\]
we usually intend 
	\[
	\theta = \hat{\theta}(x_0) \pm \sqrt{\Var\left[\hat\theta\right]}\bigg\rvert_{\theta=\hat{\theta}}	
	\]
where we explicitly write that $\hat{\theta}$ depends on the observations $x_0$ and that $\Var\left[\hat\theta\right]$ is a function of $\theta$ evaluated at $\hat{\theta}$.\todo{There are some questions left as exercises here if we want to include them, it's Saturday and I don't feel like thinking about them right now...}


\begin{example}
	One could think that having a consistent estimator means that the variance must go to zero:
	\[
	\hat{\theta}_N \text{ consistent} \implies \lim_{N \to \infty} \text{Var}\left[\hat{\theta}_N\right] = 0
	\]
	However, it is easy to find a counterexample. Consider an estimator $\hat{\theta}_N$ for the parameter $\theta = 0$ defined by the following probability mass function:
	\[
	P(\hat{\theta}_N = x) = 
	\begin{cases} 
		\frac{1}{2N} & \text{if } x = N^2 \\
		1 - \frac{1}{N} & \text{if } x = 0 \\
		\frac{1}{2N} & \text{if } x = -N^2 \\
		0 & \text{otherwise}
	\end{cases}
	\]
	
	\begin{figure}[H]
		\centering
		\input{images/consistency_counterexample.pgf}
		\caption{Should we add a caption and a label to this image??}
	\end{figure}
	
	An estimator is consistent if it converges in probability to the true parameter: 
	\begin{align*}
		\forall \varepsilon > 0: P(|\hat{\theta}_N - 0| > \varepsilon) &= P(\hat{\theta}_N = N^2) + P(\hat{\theta}_N = -N^2) \\
		&= \frac{1}{2N} + \frac{1}{2N} = \frac{1}{N}
	\end{align*}
	Taking the limit:
	\[
	\lim_{N \to \infty} P(|\hat{\theta}_N| > \varepsilon) = \lim_{N \to \infty} \frac{1}{N} = 0
	\]
	This proves that $\hat{\theta}_N \xrightarrow{P} 0$, so the estimator is \textbf{consistent}.
	
	The expected value is:
	\[
	E[\hat{\theta}_N] = (-N^2)\frac{1}{2N} + (0)(1-\frac{1}{N}) + (N^2)\frac{1}{2N} = 0
	\]
	The variance is then:
	\[
	\text{Var}(\hat{\theta}_N) = E[\hat{\theta}_N^2] = (-N^2)^2 \frac{1}{2N} + 0 \left(1 - \frac{1}{N}\right) + (N^2)^2 \frac{1}{2N} = N^3
	\]
	As $N \to \infty$, the variance $\text{Var}(\hat{\theta}_N) \to \infty$, showing that consistency does not implies variance decay.
\end{example}


In general, estimators with lower variance are naturally considered to be "better" (or at least low variance is a favorable property). However, the tale does not end here, and we define another property of point estimators: their \emph{bias}.
\begin{definition}[Bias of an estimator]
	Given an estimator $\hat{\theta}(x)$ of $\theta$ we define the \emph{bias function} as:
	\[
	b_\theta\left[\hat\theta\right]=\E\left[\hat{\theta}\right]-\theta
	\]
\end{definition}

While we talk about the bias of an \emph{estimator} $\hat\theta$, it's important to remember the dependency on the \emph{parameter} $\theta$. An unbiased estimator is one that has identically null bias for \emph{all} values of $\theta$\sidenote{For example, the sample mean is always unbiased (LLN) while a constant value has zero bias if you "guess" the parameter but biased in all other cases, and is therefore biased.}, and if this condition cannot be met for all $\theta$ it's at least favorable that the bias be much smaller than the variance. When we write
	\[
	\theta = \hat{\theta}(x_0) \pm \sqrt{\Var\left[\hat{\theta}\right]}
	\] 
	
we silently intend that the aforementioned condition is met:
	
	\[
	\left(b\left[\hat{\theta}\right]\ll\sqrt{\Var[\hat{\theta}]}\right)
	\]

We can in general try to reduce the bias of an estimator "artificially": for example we can think to banally subtract it from the estimator itself to obtain a new one. In this case, we will have:

	\[
	\hat\theta^\prime = \hat\theta - b_\theta\left[\hat\theta\right]
	\]

We note that $\hat\theta^\prime$ is only a statistic if $\frac{\partial b}{\partial \theta} = 0$, as the bias of $\hat\theta$ can depend on $\theta$ itself (not constant bias). Well, we can go one step deeper and find an \emph{estimator for the bias}, and newly define:

	\[
	\hat\theta^\prime = \hat\theta - \hat{b} 
	\]

What exactly is $\hat{b}$? We can, like we usually do for estimators, build it \emph{ad hoc} using experience to find a good estimator, but two possible candidates are $\hat{b}_\theta\left[\hat\theta\right]_{\theta=\hat{\theta}}$ or even $\hat{b}_\theta\left[\hat\theta\right]_{\theta=\hat{\theta}^\prime}$, where essentially we get rid of the problem of having $\theta$ by explicitly estimating it with different types of estimators (the example that follows should make this clearer).
The new bias will be:

	\[
	b_\theta^\prime\left[\hat\theta^\prime\right] = \E\left[\hat\theta^\prime\right] - \theta = \E\left[\hat\theta\right] - \E\left[\hat{b}\right] - \theta = b_\theta(\hat\theta) - \E\left[\hat{b}\right]
	\]

The new estimator is unbiased when the last difference is equal to zero:
	
	\[
	\E\left[\hat{b}\right] - b_\theta\left[\hat\theta\right] = 0 
	\]

If you squint hard enough at this expression you can convince yourself that it means nothing more than \emph{the estimator for the bias is unbiased itself}.

Unfortunately, as things often are, reducing the bias of an estimator in this way is often matched with increasing its variance (the explanation for this will be clearer in the next section).

\begin{example}
	To try to keep things as clear as possible, let's formulate a quick example: we have built an estimator $\hat\theta$ for $\theta$ and found that its bias is
	
	\[
	b_\theta \left[ \hat\theta \right] = \E\left[\hat\theta\right] - \theta = c \cdot \theta
	\]
	
	where $c$ is a real-valued nonzero constant. Inverting we also find that $\E\left[\hat\theta\right] = (c+1)\theta$ which will be useful later.
	The estimator is obviously not unbiased so we try to correct this by building a new estimator $\hat\theta^\prime = \hat\theta - c\theta$. Well, we don't know what $\theta$ is, that's what we are trying to estimate in the first place, and so the best we can do is to replace it with $\hat\theta$. Therefore, I can build my new estimator to be:
	
	\[
	\hat\theta^\prime = \hat\theta - c\hat\theta
	\]
	
	We recognize the second term to be what we previously called the estimator for the bias $\hat{b} = b_\theta\left[\hat\theta\right]_{\theta=\hat\theta}$. 
	
	Let's try to calculate the new bias of this estimator $\hat\theta^\prime$:
	
	\begin{align*}
	b_\theta\left[\hat\theta^\prime\right] &= \E\left[\hat\theta^\prime\right] - \theta = \E\left[\hat\theta\right] - \E\left[c\hat\theta\right] - \theta \\
	&= b_\theta\left[\hat\theta\right] - c\E\left[\hat\theta\right] = c\theta - c(1+c)\theta = -c^2\theta
	\end{align*}
	
	That's not very promising. Let's try the other suggested case, where we build the estimator as:
	
	\[
	\hat\theta^\prime = \hat\theta - c\hat\theta^\prime
	\]
	
	Where we are instead using $\hat{b} = b_\theta\left[\hat\theta\right]_{\theta=\hat\theta^\prime}$ as an estimator for the bias. We solve this implicit function to get
	
	\[
	\hat\theta^\prime = \frac{\hat\theta}{c+1}
	\]
	
	and calculating the bias once again:
	
	\begin{align*}
	b_\theta\left[\hat\theta^\prime\right] &= \E\left[\hat\theta^\prime\right] - \theta = \frac{\E\left[\hat\theta\right]}{c+1} - \theta \\
	&= \frac{c+1}{c+1}\theta - \theta = 0
	\end{align*}
	
	Which, as we have verified, is unbiased.
\end{example}
\todo{Do we wanna find the variances of these estimators too?}

When we can't make do with an unbiased estimator, we can at least try to find one that is asymptotically so:

	\[
	\forall \theta:\, \lim_{N\to\infty}b_N\left[\hat\theta\right]=0
	\]
	
\todo{More questions and counterexamples.}

\subsection{Cramer-Rao}
\label{subsubsec:CR}

Question: can we arbitrarily reduce the variance and bias of an estimator or is there a limit? Unfortunately, the answer is that there is, and we introduce an important theorem for this course:

\begin{theorem}[Cramer-Rao]
	Let's suppose that the Fisher score \ref{def:Fisher_S} exists $\forall x$ and that $\frac{\partial}{\partial\theta}$ commutes with $\E$. Considering $\hat\theta$, an estimator for $\theta$, then there exists a lower limit to the variance of $\hat\theta$:
	\[
	\Var\left[\hat\theta\right]\geq \left(1+\frac{\partial b}{\partial\theta}\right)^2 I_{\hat\theta}^{-1} \geq \left(1+\frac{\partial b}{\partial\theta}\right)^2 I_{x}^{-1}
	\]
\label{def:cramer-rao}
\end{theorem}

\begin{example}[Proof of the C-R theorem]
	We start from the following identity:
	\[
	\Var\left[\hat\theta\right]\cdot I_{\hat\theta} = \E\left[\left(\hat\theta-\E\left[\hat\theta\right]\right)^2\right]\cdot\E\left[\left(\frac{\partial\ln p(\hat\theta;\theta)}{\partial\theta}\right)^2\right]
	\]
	Where $I_{\hat\theta}$ is the Fisher Information contained in the estimator $\hat\theta$ (which we remember to be, \emph{in primis}, a statistic).
	
	From the Schwartz inequality: 
	\[
	\Var\left[\hat\theta\right]\cdot I_{\hat\theta}\geq \E\left[\left(\hat\theta-\E\left[\hat\theta\right]\right)\cdot\frac{\partial\ln p(\hat\theta;\theta)}{\partial\theta}\right]^2
	\]
	
	We'll split $\E\left[\left(\hat\theta-\E\left[\hat\theta\right]\right)\cdot\frac{\partial\ln p(\hat\theta;\theta)}{\partial\theta}\right]$ into two parts:
	
	\begin{align*}
		1.) \quad \E\left[\hat\theta\cdot\frac{\partial\ln p(\hat\theta;\theta)}{\partial\theta}\right] &= \int \hat\theta\cdot\frac{\partial\ln p(\hat\theta;\theta)}{\partial\theta} p(\hat\theta;\theta)\dd \hat\theta \\
		&= \int \hat\theta\cdot\frac{\partial p(\hat\theta;\theta)}{\partial\theta}\dd \hat\theta = \frac{\partial}{\partial\theta}\int \hat\theta\cdot p(\hat\theta;\theta)\dd \hat\theta \\
		&=\frac{\partial}{\partial\theta}\E\left[\hat\theta\right] = \frac{\partial}{\partial\theta}\left[\theta+b_\theta\left[\hat\theta\right]\right]\\
		&= 1 + \frac{\partial b}{\partial\theta}
	\end{align*}
	
	\begin{align*}
		2.) \quad \E\left[-\E\left[\hat\theta\right]\cdot\frac{\partial\ln p(\hat\theta;\theta)}{\partial\theta}\right] &= -\E\left[\hat\theta\right]\int \frac{\partial\ln p(\hat\theta;\theta)}{\partial\theta} p(\hat\theta;\theta)\dd \hat\theta \\
		&= -\E\left[\hat\theta\right]\int\frac{\partial p(\hat\theta;\theta)}{\partial\theta}\dd \hat\theta \\
		&= -\E\left[\hat\theta\right]\frac{\partial}{\partial\theta}\int p(\hat\theta;\theta)\dd \hat\theta \\
		&=-\E\left[\hat\theta\right]\frac{\partial}{\partial\theta}1 = 0
	\end{align*}
	
	Putting everything together we get the Cramer-Rao inequality:
	
	\[
	\Var\left[\hat\theta\right]\geq \left(1+\frac{\partial b}{\partial\theta}\right)^2 I_{\hat\theta}^{-1} \geq \left(1+\frac{\partial b}{\partial\theta}\right)^2 I_{x}^{-1}
	\]
	
	\qed
	
	The first equality is verified when the estimator has the minimum variance out of all the possible ones we can build from the same statistic. This is verified when we have the equality in the Schwartz inequality (when the "vectors" are "parallel"), which means when
	
	\begin{align*}
	\left(\hat\theta - \E\left[\hat\theta\right]\right)\, &||\, \frac{\partial\ln p (\hat\theta;\theta)}{\partial\theta} \Rightarrow \frac{\partial\ln p(\hat\theta;\theta)}{\partial\theta} = f(\theta) \cdot (\hat\theta-\theta-b(\theta))\\
	&\Rightarrow \ln p(\hat\theta) = \hat\theta\int f(\theta)\dd \theta - \int \left(\theta+b(\theta)\right)\cdot f(\theta)\dd \theta +c\\
	&\Rightarrow p(\hat\theta;\theta) = \exp\left[\alpha(\theta)\hat\theta+\beta(\theta)+c(\hat\theta)\right]
	\end{align*}
	
	Which means that to satisfy the first inequality we must be able to express the pdf in the exponential form written above which means that statistics that satisfy Darmois \ref{th:darmois} have minimum variance. We also note that this statistic is unbiased for $\E\left[\hat\theta\right]$, and not necessarily unbiased for $\theta$. Unfortunately, to remove the bias in $\theta$ we will have to sacrifice the minimum variance. 
	
	The other inequality is satisfied for sufficient statistics. 
\label{def:CR}
\end{example}\todo{Not all of these calculations are totally clear to me I just wrote this on the plane without trying to understand it. Therefore, this section could do with some love. Also note that I am not particularly certain of what I have written...}

This is awesome, we have an elegant expression that ties together so many concepts that we have studied and justifies their use in the theory of Probability. We also have a way to write the \textbf{Minimum Variance Bound} (MVB), expressing the limit under which the variance of an estimator cannot drop. We can also state that there certainly exists at the most one function of the parameters that can be optimally estimated and without bias\todo{Why? I don't get it. When we understand this, replace this todo with a sidenote saying "this principle goes under the name of the Lehmann-Scheffé theorem"}


 We can quantitatively define the efficiency of an estimator as
 
	\[
	\varepsilon_{\hat\theta}(\theta) = \left(1+\frac{\partial b}{\partial \theta}\right)^2 I_x^{-1}\Var^{-1}\left[\hat\theta\right]
	\]

which is defined between zero and one, and when maximal allows us to say that the estimator is \emph{efficient}.

\begin{corollary}
	If the Fisher Information is linear in $N$, then
	
	\[
	\Var\left[\hat\theta(x)\right] \geq \frac{k}{N}
	\]
	
	As the measurements increase in number, the variance decreases as $\frac{1}{N}$, though it only holds if the Likelihood satisfies some basic regularity conditions.
\end{corollary}\todo{is there a counter-example?}

Generalizing the Cramer-Rao inequality to multiple statistics and parameters:

	\[
	\Cov\left[\hat\theta\right]_{ij}\geq\left(1+\frac{\partial b_i}{\partial \theta_k}\right)I^{-1}_{kl}\left(1+\frac{\partial b_l}{\partial \theta_j}\right)
	\]
	
Where by "$\geq$" we intend that the difference between the left and right hand sides returns a positively-defined matrix.

\todo{There are some Cramer-Rao examples (and NON Cramer-Rao) here, let's work them out or should this be stuff to put in the appendix? I vote to prove here and copy results in the appendix. You decide and do, I agree with whatever.}

\subsection{Point estimates with moments}

What do we do in the unfortunate case that we cannot express our pdf as an exponential function? Well then, we don't have the \emph{MV} guaranteed, nor a sufficient statistic. How do we find a decent estimator for $\theta$ in the most general case? Well, we have one estimator for free in general, and the LLN comes to the rescue. Recall that we have stated that the arithmetic mean is a consistent estimator for $\E[x]$, and so in the worst case we can just use the good old arithmetic mean if the expected value is the parameter we are trying to estimate (as long as the latter exists and $x$ has finite variance). 

And in the other cases? Let's explore a general and powerful technique, recalling the \emph{moments} that we have talked about in section \ref{subsec:moments}. For our point estimation problem, it's natural to consider the statistics of the following kind:

	\[
	m_i = \frac{1}{N} \sum_{k=1}^{N}x_k^i
	\]
	
with, if you recall,

	\[
	m_0 = 1,\quad m_1=\bar{x}, \quad m_2 = \bar{x^2}
	\]
	
It's not hard to imagine why we'd call these the \emph{sample moments}, and they satisfy the following relation (by definition of the moments of a distribution centered at $0$):

	\[
	\E\left[m_i\right] = \mu_i^\prime
	\]
	\todo{We should rework this whole subsection}

It's not guaranteed that these exist and are finite, but when they do they prove to be helpful statistics that we can readily calculate and update as we gather new data.\todo{More junk here that I happily skipped on my first few takes of this exam} If enough moments exist, this method allows us to easily build consistent estimators with known covariance\todo{if we talk of the cov then we have to include the "junk"} and (for the central limit theorem) asymptotically normal. They are not necessarily unbiased\todo{Said in Italian: potrei aver detto una minchiata.}.

The cool thing about these moments is that they often exist finite and are easy to calculate, both in terms of computational power and weight. When you need to make some rapid decisions,\sidenote{Like, LHC GHz level of rapid} they can be pivotal to calculate "on the fly". 

Of course, every good property must be met with an unfavorable one (in certain fields called the "no free lunch principle"), and this method does not grant us particularly efficient estimators, and to increase their efficiency we naturally must calculate more, or increase the complexity of the model, which is in direct contradiction with the whole reason we chose the moment method in the first place.\todo{Bleh. Let's re-write this sentence to flow better.}

\subsection{Implicit estimators}
\label{subsec:Implicits}
Let's start with a basic fact: the expected value of a statistic will be, in general, at most a function of the parameters: $\E[\hat\theta(x)]  = f(\theta_0)$. By simply re-writing our expression

	\[
	\E[\hat\theta(x)]  = f(\theta_0) \Rightarrow \E\left[\hat\theta(x)-f(\theta_0)\right] = 0
	\]

we can decide to generalize this concept by considering a generic $F(x,\theta)$ for which $\E\left[F(x,\theta=\theta_0)\right]=0$. In other words this condition is verified only when $\theta$ assumes the "real" value $\theta_0$ from which $x$ is extracted according to $p(x;\theta_0)$. We note that this $F$ isn't a statistic, and at most can be seen as a family of statistics for each fixed $\theta$. 

For a generic $\theta$, we suppose that there exists finite (at least in a neighborhood of $\theta_0$):

	\[
	\E\left[F(x,\theta)\right]\equiv h(\theta)\quad \text{where}\quad h(\theta_0)=0
	\]

We now define $k_N(\vec{x},\theta) = \frac{1}{N}\sum_{i=1}^N F(x_i,\theta)$. We will have that

	\[
	k_N(\vec{x},\theta) \xrightarrow[N \to \infty]{P} \E\left[F(x,\theta)\right] = h(\theta)
	\]

Where we have indicated with a $P$ the convergence in probability for high statistics. 


\begin{definition}\label{def:implicit_estimator}[Implicit estimator]
	
	Let's now introduce what we want to talk about in this section: the (family of) implicit estimators $\hat\theta_N$ and define them as the solution to the equation: 
	
	\[
	k_N(\vec{x}, \hat\theta) = 0
	\]
	
 	If in a neighborhood of $\theta_0$:
 	
 	\begin{itemize}
 		\item $k_N(\vec{x},\theta)\in\mathbb{C}^1$
 		\item $\exists\E\left[\frac{\partial k_N(\vec{x},\theta)}{\partial \theta}\right]$
 		\item $\E\left[\frac{\partial k_N(\vec{x},\theta)}{\partial\theta}\right]_{\theta=\theta_0}\neq 0$
 	\end{itemize}
 	
 	then it can be shown that $\hat\theta_N$ is, for at least one of the solutions, a consistent estimator for $\theta$:
 	
	 	\[
	 	\hat\theta_N \xrightarrow[N \to \infty]{P} \theta_0
	 	\]
	 	
\end{definition}

Let's take a step back and look at what we are dealing with. Before this section we talked about (without really giving them a name) \emph{explicit estimators}, where we have an expression of the type $\hat\theta = s(x)$, so by simply measuring some data and doing an operation on it (like taking the sample mean) we directly (explicitly) obtain a consistent estimator for the unknown parameter $\theta$. Great! Implicit estimators, instead, are slightly different: we have an expression of the type $k_N(\vec{x}, \hat\theta_N) = 0$. Here, the estimator $\hat\theta_N$ is defined not by a direct formula, but as the root of an estimating function. If certain conditions are met then $\hat\theta_N$ is a consistent estimator for $\theta$.

Finding the root of this function is not usually an easy task, and luckily we won't have to explicitly, in many cases. However, this formalism is important to introduce for a very important implicit estimator that we will use over the course of this text.

\section{Maximum Likelihood estimator}

Let's study the particular case of using the Fisher Score (\ref{def:Fisher_S}) as $F(x,\theta) = S_x(\theta) = \frac{\partial \ln L_x(\theta)}{\partial\theta}$. We have already seen that in non-particular cases the expected value of the Fisher score is null, and we define the function 

	\begin{align*}
		k_N(\vec{x},\theta) &= \frac{1}{N}\sum_{i=1}^N S_{x_i}(\theta) = \frac{1}{N}\sum_{i=1}^N \frac{\partial\ln L(x_i,\theta)}{\partial\theta} \\
		&= \frac{1}{N}\frac{\partial}{\partial\theta}\sum_{i=1}^N \ln L(x_i,\theta)=\frac{1}{N}\frac{\partial}{\partial\theta}\ln \left(\prod_{i=1}^{N}L(x_i,\theta)\right) \\
		&= \frac{1}{N}\frac{\partial}{\partial\theta}\ln L_{tot}(\vec{x},\theta)
	\end{align*}


We can see that
	
	\[
	\E\left[\frac{\partial k_N(\vec{x},\theta)}{\partial\theta}\right]_{\theta=\theta_0} = \E\left[\frac{1}{N}\frac{\partial^2\ln L_{tot}(\vec{x},\theta)}{\partial\theta^2}\right]_{\theta=\theta_0}= \frac{-I_F}{N} \neq 0
	\]

From this we see that the three conditions required in the definition of implicit estimator \ref{def:implicit_estimator} are satisfied. Thus we conclude that at least one of the solutions $\hat\theta$ of the Likelihood equation:

\begin{equation}
	\frac{\partial}{\partial\theta} \ln L_{tot}(\vec{x},\hat\theta_{MLE}) = 0 
	\tag{Likelihood equation}
	\label{eq:likelihood}
\end{equation}

is a consistent estimator for $\theta$. 

This estimator is the \textbf{Maximum Likelihood Estimator} (MLE).

Let's go through some of its nice properties:
\begin{itemize}
	\item The Fisher Information is defined positive, and thus the MLE is, by definition, a maximum.
	\item It respects the Likelihood Principle.
	\item It is invariant for parameter variable change (just like the Likelihood).
	\item Asymptotically, it is efficient and unbiased (goes like $\frac{1}{N}$) (\textbf{this is not guaranteed for a finite sample!})
	\item Asymptotically, it is also normal, with: \[p(\hat\theta;\theta_0) \rightarrow \N\left(\hat{\theta}; \mu=\theta_0,\sigma^2=\frac{1}{I_F}\right)\].
\end{itemize}

Since we know a thing or two about compromises, and because lunch is never free, we add that the MLE, being quite powerful, also requires minimization technique, as well as the entire dataset, to be calculated, making it much slower and computationally expensive to calculate than the moments. However, it is still a favorite option to use when possible. 


\subsection{Variance of the MLE}

Now that we understand the power of the MLE, we see why it is so widely used in real-world problems. In more complex cases, especially with large datasets, the maximum of the total likelihood $L_{\text{tot}}$ is typically found numerically using optimization algorithms (\emph{fitting})\todo{From here we start talking about \emph{fitting} but we never defined what a fit is in the first place}. The question is: how can we evaluate the variance of the MLE in these cases?

In the asymptotic limit, we can use of the Cramer-Rao theorem \ref{def:cramer-rao} to state that, if $I$ is linear:

	\[
		\Var\left[\hat\theta_{MLE}\right]\simeq\frac{1}{NI_1}
	\]
	
where $I_1$ is the Fisher information of a single observation. 
In the complex cases mentioned above, it is difficult to compute $I_1$ directly. However, we can write (under the usual regularity conditions that allow us to commute the expected value with the derivative):

	\[
		N \cdot I_1= -N\E\left[\frac{\partial^2\ln L(\theta)}{\partial\theta^2}\right]=-N\frac{\partial^2}{\partial\theta^2}\E\left[\ln L(\theta)\right]
	\]
	
If the expected value exists, it can be consistently estimated from the data using our trusty LLN:

\begin{align*}
	N \cdot I_1=-N\frac{\partial^2}{\partial\theta^2}\E\left[\ln L(\theta)\right]	&\simeq -N \frac{\partial^2}{\partial\theta^2} \frac{1}{N}\sum_{i=1}^N\ln L_i(\theta) \\
	&= -\frac{\partial^2\ln L_{tot}(\theta)}{\partial\theta^2}
\end{align*}

	Using the fact that $\hat\theta_{MLE}$ is a consistent estimator for $\theta$, we finally obtain:
	
	\[
		 \Var(\hat\theta_{MLE})\simeq-\left(\frac{\partial^2\ln L_{tot}(\hat\theta_{MLE})}{\partial\theta^2}\right)^{-1}
	\]
	
In other words: \emph{the asymptotic variance of the MLE can be evaluated using the second derivative of the total likelihood $L_{tot}$}. 
One advantage of this formula is that the total likelihood has already been computed and maximized numerically, so evaluating its second derivative is computationally inexpensive.

Generalizing to $n$ dimensions:

	\[
		\Cov\left(\hat{\theta}_i,\hat{\theta}_j\right)=-\left(\frac{\partial^2\ln L_{tot}(\hat{\theta})}{\partial\theta_i\partial\theta_j}\right)^{-1}
	\]
	
This result is widely used in fitting software (that you've probably used thousands of times) to evaluate the errors on fitted parameters!
It is important to stress that this is an asymptotic result and we invite the reader to take another look at all the approximations made throughout the derivation.

\section{Fitting histograms}
\label{sec:fitting_hist}

In the quite common case where no sufficient statistic is available, a compact representation of the data can be obtained using histograms. One should always keep in mind that the information content is reduced: $I_{\text{hist}} \leq I_x$.

Let $k_i$ be the number of events $\{x_j\}$ falling into the bin $i \in [1,n]$. Each bin is defined by a subset of the event space $S_i \subset X$. The bins can be constructed in several ways (equal width, equal probability, etc.), but this choice is not relevant for what follows.

A histogram is therefore a statistic defined by the vector of counts $\{k_i\}$, which are functions of the $N$ observed events $\{x_j\}$. To each bin we associate a probability:

	\[
		\mathfrak{p}_i(\theta)\equiv P(\vec{x}\in S_i)=\int_{S_i}p(\vec{x};\theta)dx
	\]

We now distinguish two scenarios, depending on whether $N$ is fixed or not.

\subsection{N-fixed histograms}

If $N$ is fixed \emph{a priori}, the counts $\{k_i\}$ follow a multinomial distribution\todo{Insert here the ref to the appendix with all the distribution defined}:

	\[
		P(\{k_i\};\mathfrak{p}_i(\theta),N)=\frac{N!}{\prod_i{k_i!}}\prod_i \mathfrak{p}_i^{k_i}(\theta)
	\]
	
the log-Likelihood is then:

	\[
		\ln L(\theta)=\sum_{i}k_i \ln \mathfrak{p}_i(\theta)-\sum_i\ln k_i! + \ln N!
	\]
	
Thus the Likelihood equation:

	\[
		\frac{\partial\ln L(\theta)}{\partial\theta}\bigg\rvert_{\theta=\hat\theta}=\sum_i \frac{k_i}{\mathfrak{p}_i(\theta)}\frac{\partial\mathfrak{p}_i(\theta)}{\partial\theta}\bigg\rvert_{\theta=\hat\theta}=0
	\]

We now want to compare the inference obtained from the binned data with that obtained from the unbinned (raw) data.

In the limit of "small" bins:

	\[
		\mathfrak{p}_i(\theta)\simeq V_i\cdot p(x_i;\theta)
	\]

where $V_i=\int_{S_i}d\vec{x}$ is the "volume" of the $i$-th bin and $x_i$ is the value of $x$ in the center of the $i$-th bin. The reader is warned that this approximation is sometimes swept under the rug. 

We can than write:

\begin{align*}
	\ln L (\theta) &\simeq \text{Const} + \sum_i k_i \ln (V_i p(x_i;\theta)) \\
	&= \text{Const}^\prime + \sum_i k_i \ln p(x_i;\theta)
\end{align*}

In the limit of infinitely fine binning ($n \to \infty$), each bin contains at most one event, so that $k_i \in \{0,1\}$ and:

\begin{align*}
	\ln L(\theta) &= \text{Const} + \sum_{i=1}^{N}\ln p(x_i;\theta)\\
	&= \text{Const} + \ln L_{\text{unbinned}}(\theta) 
\end{align*}

\emph{Therefore, in the limit of infinitesimally small bins, the binned and unbinned likelihoods coincide up to an additive constant. This implies that both approaches lead to the same inference asymptotically.}

\subsection{Variable-$N$ histograms}
\label{subsec:variable_N_hist}

When $N$ is not fixed (for instance, when observations are collected over a fixed time interval), it becomes itself a random variable, following a Poisson distribution:\todo{Insert here the ref to the appendix with all the distribution defined}
\[
	N \sim \mathcal{P}(N; \mu(\theta))
\]
where $\mu$ is the expected number of total events:
\[
	\mu(\theta)\equiv\E\left[N(\theta)\right]=\E\left[\sum_i k_i(\theta)\right]
\]

In this case, the bin counts $\{k_i\}$ are independent and each follows a Poisson distribution:
\begin{align*}
	P(\{k_i\};\theta)&=\prod_i\mathcal{P}(k_i;\mu_i(\theta))\\
	&=\prod_i \frac{\mu_i(\theta)^{k_i}e^{-\mu_i(\theta)}}{k_i!}
\end{align*}
where $\mu_i(\theta)$ is the expected number of events in the $i$-th bin.

If $\mathfrak{p}_i(\theta)$ denotes the probability for an event to fall into the $i$-th bin, then
\[
	\mu_i (\theta) \equiv \E [ k_i (\theta) ] = \E [\mathfrak{p}_i(\theta)N(\theta)] = \mu(\theta)\mathfrak{p}_i(\theta)
\]
This relation follows from the fact that, conditionally on $N$, the counts $\{k_i\}$ follow a multinomial distribution with probabilities $\{\mathfrak{p}_i\}$, and taking the expectation over the Poisson-distributed $N$ yields independent Poisson variables with means $\mu\,\mathfrak{p}_i$.

The log-likelihood is therefore:
\begin{align*}
	\ln L(\theta)&=\sum_i k_i \ln \mu_i(\theta)-\sum_i \mu_i(\theta)-\sum_i \ln k_i!\\
	&=\sum_ik_i\ln (\mu (\theta) \mathfrak{p}_i(\theta))-\mu(\theta) + Const\\
	&=\sum_i k_i \ln \mathfrak{p}_i(\theta) + \sum_i k_i \ln \mu (\theta) -\mu(\theta) + Const
\end{align*}

Using the same "small bins" approximation as before ($\mathfrak{p}_i(\theta)\simeq V_i\,p(x_i;\theta)$), we obtain:
\begin{align*}
	\ln L(\theta)&=\sum_i k_i \ln (V_i p(x_i;\theta)) + N \ln \mu(\theta) - \mu(\theta) + Const \\
	&=\sum_i k_i \ln p(x_i;\theta)+ N \ln \mu(\theta) - \mu(\theta) + Const'
\end{align*}

We notice that this log-likelihood differs from the multinomial case only by the additional term $N \ln \mu(\theta) - \mu(\theta)$ (up to irrelevant constants).

We must now distinguish between two scenarios, depending on whether $\mu$ depends on $\theta$ or not:

\begin{itemize}
    \item If $\frac{\partial \mu}{\partial \theta} = 0$, then $\mu$ is an independent parameter and the likelihood equations include an additional condition:
    \[
		\frac{\partial \ln L}{\partial \mu} = 0 
		\;\Rightarrow\; \frac{N}{\mu} - 1 = 0 
		\;\Rightarrow\; \hat{\mu} = N
	\]
	as expected. The remaining equations coincide with those of the multinomial case.

    \item If $\frac{\partial \mu}{\partial \theta} \neq 0$, then $\mu$ carries information about the system, and the observable $N$ contributes non-trivially to the inference of $\theta$. In this case, even in the infinitesimal binning limit, the likelihood differs from the multinomial case:
	\begin{equation}
		\ln L(\theta)=\sum_i^{N}\ln p(x_i; \theta) + N \ln \mu(\theta) - \mu(\theta) + Const
		\label{eq:extended_likelihood}
	\end{equation}
	which is known as the \textbf{extended likelihood}. It is easy to verify that, given $N$ observed events, the extended likelihood can be written as the product of the usual likelihood and a Poisson term in $N$:
	\[
		L_{ext}(\theta)=\mathcal{P}(N;\mu(\theta))L_{tot}(\theta)=\frac{\mu(\theta)^Ne^{-\mu(\theta)}}{N!}L_{tot}(\theta)
	\]
\end{itemize}

\subsection{High-statistics histograms}

One of the most common situations in real-life applications is the case of high statistics in every bin:
\[
	\forall i: \mu_i(\theta)=\E[k_i] \gg 1
\]
In this regime, the distribution of counts in each bin can be approximated by a normal distribution (see \ref{th:central_limit}). Therefore, the likelihood can be written as
\[
L(\theta) = \prod_i \exp\left(-\frac{(k_i - \mu_i(\theta))^2}{2\,\sigma_i^2(\theta)}\right) \cdot Const
\]
Taking the logarithm and dropping irrelevant constants:
\begin{align*}
	\ln L(\theta)
	&= -\frac{1}{2}\sum_i \frac{(k_i - \mu_i(\theta))^2}{\sigma_i^2(\theta)} \\
	&= -\frac{1}{2}\sum_i \frac{(k_i - \mu_i(\theta))^2}{\mu_i(\theta)}
	\;\equiv\; -\frac{1}{2}\chi^2
\end{align*}
where in the last step we used $\sigma_i^2(\theta) = \mu_i(\theta)$, as appropriate for Poisson statistics.

Thus, the log-likelihood reduces to the \textbf{Pearson chi-squared} statistic. In the high-statistics limit, minimizing $\chi^2$ is equivalent to maximizing the likelihood.

A further simplification is sometimes introduced:
\begin{equation*}
	\chi^2 = \sum_i \frac{(k_i - \mu_i(\theta))^2}{\mu_i(\theta)}
	\;\simeq\;
	\sum_i \frac{(k_i - \mu_i(\theta))^2}{k_i}
	\tag{Modified $\chi^2$}
	\label{eq:modified_chi2}
\end{equation*}
This expression is asymptotically equivalent, and has the advantage of being quadratic in $\theta$ if $\mu_i(\theta)$ is linear in $\theta$. However, it has significantly worse convergence properties, and the resulting estimator can be biased, sometimes with a bias larger than its statistical uncertainty.

\section{Robustness of a statistic}
We have seen thus far the basic principles of classic point estimators, where in most cases the pdf of the observable is known almost to perfection. However, this is not always the case, and it's important to talk about one final concept: the \emph{robustness} of a statistic. 

\begin{definition}[Robustness of a statistic]
	The \emph{robustness} measures how well a statistic maintains its properties even if the underlying distributional assumptions are incorrect. 
\end{definition}

In a sense, a robust statistic is more trustworthy, as it is less volatile to outliers, model inconsistencies, or noise. The last thing we want as experimental scientists is to use methods that break down for tiny variations of the initial hypotheses, and for this reason robustness is a sought out property.

\begin{example}
	A good estimator for the variance of a Gaussian distributed random variable is the \emph{root-mean-square} (RMS), which, being the MLE, is asymptotically efficient:
	\[
	\hat\sigma \equiv s_N = \sqrt{\frac{1}{N}\sum_i (x_i - \bar{x})^2}
	\]
	where $\bar{x}$ is the sample mean.
	
	This is not the only possible estimator for the variance. We could use, for example, the \emph{mean-absolute-deviation} (MAD):
	\[
	d_N = \frac{1}{N}\sum_i |x_i - \bar{x}| \quad \rightarrow \quad \hat\sigma' \equiv \sqrt{\frac{\pi}{2}}\, d_N
	\]
	where the multiplicative factor makes the estimator asymptotically unbiased for a Gaussian distribution.
	
	We may ask which estimator performs better. Using the concept of efficiency, one finds:
	\[
	r \equiv \lim_{N \to \infty} \frac{\varepsilon(\hat\sigma')}{\varepsilon(\hat\sigma)} = 0.876
	\]
	This means that the efficiency of $d_N$ is significantly worse than that of $s_N$. One might conclude that this settles the question. However, the story does not end here.
	
	Consider the case where we are not completely certain about the underlying pdf. For example, suppose the data are contaminated by a small fraction $f$ of observations with a larger variance $\sigma_2$.
	We model the total pdf as a mixture of two Gaussians with the same mean $\mu$ but different variances (take $\sigma_2 = 3\sigma_1$):
	\[
	p(x; \mu, \sigma_1, f) = (1-f)\mathcal{N}(x;\mu,\sigma_1) + f\mathcal{N}(x;\mu, 3\sigma_1)
	\]
	\begin{figure}[H]
		\centering
		\input{images/historical_example.pgf}
		\caption{Should we add a caption and a label to this image??}
	\end{figure}
	We consider small values of $f$ (e.g.\ $f=0.05$), so most of the data follow the distribution with variance $\sigma_1$. This does not seem like an extreme situation, only a mild contamination. Let us now look at the efficiency ratio $r$ for different values of $f$:
	\begin{center}
	\begin{tabular}{||c|c||}
		\hline
		$f$ & $r$ \\
		\hline
		0 & 0.876 \\
		\hline
		0.001 & 0.950\\
		\hline
		0.002 & 1.02\\
		\hline
		0.050 & \textbf{2.04}\\
		\hline
	\end{tabular}
	\end{center}
	The result is remarkable: with a contamination of $f=0.05$, $d_N$ becomes about twice as efficient as $s_N$! The theory we have built so far is fragile.
	This simple example, with historical relevance, highlights the importance of robustness as a desirable property of an estimator.
\end{example}\todo{I won't talk about the history of the topic and the Eddington-Fisher debate without proper sources. We could add the historical side of things later.}

As we can see, the theory discussed up until now is rather fragile, and we'd like to find some robust estimators that we can trust even when we don't have full knowledge of the distribution of the data. 

\subsection{Least Squares method}
Let's talk about a robust estimator that is consistent even when the underlying pdf is not known. We begin by remembering what we talked about in section \ref{subsec:Implicits} about implicit estimators: all we need is a general function $F$ such that $\E\left[F(x,\theta)\right] = 0$ to build one, and we shall consider the function $F\left(k_i,\theta\right) = k_i -\mu_i(\theta) = k_i - \E\left[k_i\right]$. It's clear to see that the expected value of this function is indeed equal to zero, and we wish to transform it while maintaining this property.\todo{Is it clear here that what comes is not limited to histograms but it can be used every time there's a set of observation depending on a parameter? I think it is -J}

We apply a generic linear transformation:
	\[
	F_j\left(k_i,\theta\right) = \sum_i V_{ij}\left(k_i - \mu_i(\theta)\right)
	\]

And multiply by an arbitrary vector:
	\[
	F\left(k_i,\theta\right) = \sum_{i,j} \frac{\partial \mu_j(\theta)}{\partial\theta}V_{ij}\left(k_i - \mu_i(\theta)\right)
	\]

Let's verify the conditions for the existence of at least one consistent estimator through the solution of the above equation. We need to check that the expected value of its derivative exists not null, as described in definition (\ref{def:implicit_estimator}).
\begin{align*}
	\frac{\partial F\left(k,\theta\right)}{\partial\theta} &= -\sum_{i,j}\frac{\partial \mu_i(\theta)}{\partial\theta}V_{ij}\frac{\partial \mu_j(\theta)}{\partial\theta} + \cancelto{0}{\sum_{i,j}\left(k_j - \mu_j(\theta)\right)V_{ij}\frac{\partial^2 \mu_j(\theta)}{\partial\theta^2}}\\
	&\rightarrow -\sum_{i,j}\frac{\partial \mu_i(\theta)}{\partial\theta}V_{ij}\frac{\partial \mu_j(\theta)}{\partial\theta}
\end{align*}

This expression is different from zero if $V$ is positively-defined (as does its expected value), in which case at least one consistent estimator for $\theta$ exists. 

Looking back at the equation we want to solve: 

	\[
	\sum_{i,j} \frac{\partial \mu_j(\theta)}{\partial\theta}V_{ij}\left(k_i - \mu_i(\theta)\right) = -\frac{1}{2}\frac{\partial}{\partial\theta}\sum_{i,j} \left(k_j - \mu_j(\theta)\right)V_{ij}\left(k_i - \mu_i(\theta)\right) = 0
	\]

Which corresponds to minimizing a quadratic form. In the case of $V_{ij}$ equal to the identity\sidenote{We could also easily define it as diagonal with positive elements and obtain a weighted average.}, we have
	\[
	F(k_i,\theta) = - \frac{1}{2}\frac{\partial}{\partial\theta}\left(k_i-\mu_i(\theta)\right)^2
	\]

Whose roots we recognize as the estimators that \emph{minimize the sum of squares}. This \textbf{Least Squares Method} (LSM) allows us to find consistent estimators for the parameters $\theta$. 

Coming from the implicit estimator formalism, it follows that the consistency does not depend from the choice of $V$ (granted it is positively-defined) nor from the pdf, making it a way to find robust estimators. However it is necessary that $\exists\; \mu_i(\theta)\equiv \E\left[k_i\right]$, which could be seen as a potential weakness for distributions where this is not the case, for example ones with heavy tails or many outliers.

\begin{example}
	Let's look at one specific, yet didactic, case: where the expected values are linear functions of the parameters $\mu_i = A_{ij}\theta_j$.
	
	We minimize the squares:
	
	\begin{align*}
		0 &= \sum_{i,j} \frac{\partial \mu_j(\theta)}{\partial\theta}V_{ij}\left(k_i - \mu_i(\theta)\right) = \sum_{i,j} A_{jk} V_{ij}\left(k_i - A_{il}\theta_l\right) \\
		&\Rightarrow \sum_{i,j} A_{jk} V_{ij}k_i = \sum_{i,j} A_{jk} V_{ij}A_{il}\theta_l \\
		&\Rightarrow A^T V\cdot \vec{k} = A^T VA\cdot\vec{\theta}\Rightarrow \hat{\vec{\theta}} = \left(A^T VA\right)^{-1}\left(A^TV\right)\cdot\vec{k}
	\end{align*}
	
	This estimator is not only consistent but also unbiased:
	
	\begin{align*}
		b_{\vec\theta}\left[\hat{\vec{\theta}}\right] &= \E\left[\hat{\vec{\theta}}\right] - \vec\theta = \E\left[A^{-1}V^{-1}\left(A^T\right)^{-1}A^TV\cdot \vec{k}\right] - \vec\theta \\
		&= A^{-1}\E\left[\vec{k}\right] - \vec\theta =  A^{-1}\vec{\mu}(\theta) - \vec\theta &\\
		&=A^{-1}A\vec{\theta} - \vec\theta = 0 
	\end{align*}
	
	We add one last property to this, and in the case where $V_{ij}^{-1} = \text{cov}\left[k_i,k_j\right]$ then the estimator above is the one with minimum variance among all unbiased estimators that are linear functions of the data. We call this the \textbf{B.L.U.E} (the \emph{Best Linear Unbiased Estimator}). If the pdf of $k_i$ is a Normal distribution, then this estimator is also the MVB. 
	
	This is the Gauss-Markov theorem, which we won't prove nor go into further detail upon. 
\end{example}

Let's study a more generalized version of the least squares method:
\begin{definition}[Norm minimization]
Given a set of N observations $\{x_i\}$ of a real observable $X$, we consider the quantity (depending on parameter $p$)
	
	\[
	L_p(m) = \sum_i^N \left|x_i - m\right|^p
	\]
	
We build an estimator $\hat{m}_p=\arg\min_m L_p(m)$ that estimates a certain "location parameter" $m$ consistently and without bias. Changing $p$ changes what meaning we attribute to the parameter $m$ that we are estimating.
\end{definition}

We note that $\hat{m}$ is a consistent estimator for the mean $\mu$, and in the case of a symmetric distribution it is also unbiased.

As we stated, for different values of $p$ we are estimating different quantities:
\begin{itemize}
	\item When $p\rightarrow-\infty$ $m$ represents the \emph{mode}, which signifies the single value that appears most frequently. We are minimizing the amount of values that do not equal $m$
	\item When $p=1$ $m$ represents the \emph{median}\sidenote{We will familiarize ourselves with the median in the next section.}, which is the value where half the data lies above and below it. We are minimizing the sum of absolute differences.
	\item When $p=2$ $m=\mu$, and we are estimating the mean as in the standard least squares method. 
	\item When $p\rightarrow\infty$ $m$ represents the mid-range, which is the value exactly halfway between the maximum and minimum value in the dataset.
\end{itemize}

Depending on the type of data we have, and our interests, we can choose $p$ and use this method to find a consistent estimator for the location parameter we have chosen.

\subsection{A robust estimator: the median}
\begin{definition}[Median]
	The \emph{median} $m_1$ is the value for which there is an equal probability of being greater or less than: $P(x > m_1) = P(x<m_1) = \frac{1}{2}$. In other words, it is the value for which the cumulative density function is one half. 
\end{definition}

The sample median 
	\[
	\hat{m}_1 = \arg\min_m\left[\sum_i^N\left|x_i-m\right|\right]
	\]
provides a consistent estimate for the median, as we have seen in the previous section.

We note that the median always\sidenote{except for pathological cases involving Dirac deltas...} exists, even when $\E[x]$ does not, which makes it a robust estimate for the expected value. Also, it is invariant for monotonous transformations of the observable, and basically as long as order is conserved so is the median. 


\begin{figure}[H]
	\centering
	\input{images/median_invariance.pgf}
	\caption{Invariance of the median for the monotonous transformation $\exp(x)$. Since order is conserved, the fourth point remains the median of the transformed data.}
\end{figure}


\begin{theorem}
	If $p(x)$ is symmetric, and $\mu = E[x]$ exists, then the median is equal to the mean $m_1 = \mu$
	
	If $p(x)$ is not symmetric but $\Var[x]$ exists finite, then 
	\[
	\left|m_1-\mu\right| \leq \sigma
	\]
	which, put into words, means that the median is less than $1\sigma$ away from the mean. 
\end{theorem}
\todo{We could split these into two and prove them individually or not prove, idrk}

The sample median is such an important robust estimator that we also wish to see how it is distributed, since this will be independent of the pdf from which the data was extracted.

Assume that $p(x)$ is a continuous pdf with a unique median $m$. Let us draw an i.i.d. sample of size $2k+1$ and consider the central order statistic, i.e. the sample median $\hat m_1$.

To derive its distribution, fix a generic value $x$ and classify the observations into three bins:
\[
\begin{cases}
	x_i < x, \\
	x_i \in [x, x+dx], \\
	x_i > x.
\end{cases}
\]

The corresponding probabilities are:
\[
\begin{cases}
	P(x_i < x) = F(x), \\
	P(x_i \in [x,x+dx]) = p(x)\,dx, \\
	P(x_i > x) = 1 - F(x).
\end{cases}
\]

For $x$ to coincide with the sample median, exactly $k$ observations must fall below $x$, $k$ above $x$, and one in the infinitesimal interval $[x,x+dx]$.

This leads to a multinomial term:
\[
P =
\frac{(2k+1)!}{k!1!k!}\,[F(x)]^k [1-F(x)]^k \, p(x)\,dx.
\]

Dividing by $dx$, we obtain the probability density function of the sample median:
\begin{equation}
	p_1(\hat m_1, m_1)
	=
	\frac{(2k+1)!}{(k!)^2}
	\, p(\hat m_1)\,
	[F(\hat m_1)]^k\,[1-F(\hat m_1)]^k.
	\label{eq:median_pdf}
\end{equation}

Let's study the asymptotic case of high statistics:
	\[
	\ln p_1\left(\hat{m}_1;m_1\right) = \ln\left(\left(2k+1\right)!\right) -2\ln (k!) +k\ln\left[F\left(\hat{m}_1\right)\left(1-F\left(\hat{m}_1\right)\right)\right] + \ln p\left(\hat{m}_1\right)
	\]

Approximating with Stirling and keeping only the leading terms:

\begin{align*}
	\ln p_1\left(\hat{m}_1;m_1\right) &\sim \left(2k+1\right)\ln\left(2k+1\right) - 2k\ln k + k\ln\left[F\left(\hat{m}_1\right)\left(1-F\left(\hat{m}_1\right)\right)\right] \\
	&= 2k\ln\left[\frac{2k+1}{k}\right]+\ln\left(2k+1\right) + k\ln\left[F\left(\hat{m}_1\right)\left(1-F\left(\hat{m}_1\right)\right)\right] \\
	&\rightarrow k\ln 4 + \ln\left(2k+1\right) + k\ln\left[F\left(\hat{m}_1\right)\left(1-F\left(\hat{m}_1\right)\right)\right] \\
	&= k\ln\left[4F\left(\hat{m}_1\right)\left(1-F\left(\hat{m}_1\right)\right)\right] + \ln\left(2k+1\right) \\
	&\Rightarrow p_1\left(\hat{m}_1;m_1\right) \simeq \left(2k+1\right)\exp\left(k\ln\left[4F\left(\hat{m}_1\right)\left(1-F\left(\hat{m}_1\right)\right)\right]\right)
\end{align*}

We observe that the argument of the logarithm is always smaller than 1, except when $\hat{m}_1 = m_1$. Therefore:
\[
\lim_{k \to \infty} p_1(\hat{m}_1; m_1) = 0 \quad \forall \hat{m}_1 \neq m_1
\]

This implies that the distribution of $\hat{m}_1$ converges (in the distributional sense) to a Dirac delta centered at $m_1$:
\[
p_1(\hat{m}_1; m_1) \;\longrightarrow\; \delta(\hat{m}_1 - m_1)
\]

and thus $\hat{m}_1$ is a consistent estimator of $m_1$.

If we now go back to the asymptotic expression of the pdf:
\[
p_1(\hat{m}_1; m_1) \simeq (2k+1)\exp\left(k \ln\left[4F(\hat{m}_1)\left(1 - F(\hat{m}_1)\right)\right]\right)
\]

we can perform a Taylor expansion of the cumulative distribution around $\hat{m}_1 = m_1 + \varepsilon$:
\[
F(\hat{m}_1) \simeq F(m_1) + p(m_1)\varepsilon = \frac{1}{2} + p(m_1)\varepsilon
\]
where we used: $\frac{\partial F(m_1)}{\partial m_1}=p(m_1)$.

Substituting into the previous expression:
\begin{align*}
	p_1(\hat{m}_1; m_1)
	&\propto \exp\left(k \ln\left[4\left(\frac{1}{2} + p(m_1)\varepsilon\right)\left(\frac{1}{2} - p(m_1)\varepsilon\right)\right]\right) \\
	&= \exp\left(k \ln\left(1 - 4p^2(m_1)\varepsilon^2\right)\right) \\
	&\simeq \exp\left(-4k p^2(m_1)\varepsilon^2\right)
\end{align*}

we remember that $\varepsilon = \hat m_1 - m_1$, thus:
\[
p_1(\hat{m}_1; m_1)\simeq \exp\left(-4k p^2(m_1)(\hat m_1 - m_1)^2\right)
\]

which can be rewritten as:
\[
p_1(\hat{m}_1; m_1)\simeq \exp\left(-\frac{(\hat m_1 - m_1)^2}{2\sigma^2}\right)
\]

where\sidenote{we used $N = 2k+1 \simeq 2k$}: 
\[
\sigma^2 = \frac{1}{8k\,p^2(m_1)} \simeq \frac{1}{4N\,p^2(m_1)}
\]
showing that the distribution is asymptotically Gaussian in a neighborhood of $m_1$.

Therefore, the sample median $\hat{m}_1$ is:
\begin{itemize}
	\item consistent
	\item asymptotically normal
	\item with variance:
	\[
	\Var[\hat{m}_1] = \frac{1}{4N\,p^2(m_1)}
	\]
\end{itemize}

This shows that $\hat{m}_1$ has a finite, non-zero asymptotic efficiency.

We enjoy all of these properties independently of the fact that the variance of the observable exists finite or not (and therefore, the sample median is used in cases like the Cauchy distribution). With this we hope to have convinced the readers of just how robust the sample median is as an estimator for the median.

\todo{Hodges-Lehmann? I am feeling lazy but if it ever pops up we could end up including it...}

\todo{Other considerations that I feel like can be skipped or added later, mode and midrange I already touched upon, and truncated means are meh...}

\subsection{Truncated means}

If we wish to retain the familiar sample mean while improving its robustness, we can introduce truncation. Given N observations, we discard n values from each tail, leaving a central fraction of the data. The retained fraction is:
\[
	r= \frac{N-2n}{2N}
\]
This procedure significantly reduces sensitivity to heavy tails, making it useful for distributions such as the Cauchy or Landau. Note that for $r=1/2$ ($n=0$) we recover the standard arithmetic mean, while in the maximal trimming case $r=0$ ($n=N/2$) we recognize the median. Choosing the optimal trimming fraction typically boils down to solving a min/max problem. On one hand, removing extreme observations can improve efficiency by focusing on higher-quality (less noisy) data. On the other hand, efficiency decreases as the effective sample size is reduced. The goal is therefore to find the sweet spot between robustness and statistical efficiency, depending on the specific application.

\section{Systematic uncertainty}

Because an unknown pdf compromises our statistical certainty, we must adopt robust principles to minimize information loss. This specific form of ambiguity, which stems from a partial understanding of the pdf, is what we term \textbf{systematic uncertainty}. Under this definition, we can see robustness as the property of not being sensible to systematic uncertainties. 

Let's assume we are measuring something that has independent counts, and so we assume a Poisson distribution, where the best estimator\sidenote{The calculation for the MLE for the $N$ Poisson extractions is found in \textcolor{orange}{REF}} for $\mu$ is $\hat\mu = \frac{\sum_i k_i}{N}$ 

Now, let's assume that each measurement is affected by some background $b$, of we don't know the distribution. It's simple to see that $\hat\mu$ now becomes is $\hat\mu = \frac{\sum_i k_i}{N}-b$, which we don't fully know. We can't hope to estimate $b$ without any external information (remember, for a point estimate we need some data to build a statistic), and our original $k_i$ measurements have $b$ and $\mu$ paired together, so we can't simultaneously estimate them independently. In a certain sense, we don't even care about how $b$ is distributed, as it is just an annoying background that is not the main physics we wish to study (appropriately called "nuisance parameter"). This means that $\mu$ is afflicted by systematic uncertainty. 

Usually, I'll have some information on $b$, like restricting it within a range:

	\[
	b_0 - \frac{\delta}{2} \leq b \leq b_0 + \frac{\delta}{2}
	\]

This type of information can come from physical limitations, previous predictions, experimental data (like a "dark" measurement), and interval estimations (which we will talk about in the next chapter). 

In any case, we say \emph{something} about $b$, and, for example, we can expand $\hat\mu$ to include this:

	\[
	\hat\mu = \left(\frac{\sum_i k_i}{N}-b_0\right)\pm \sqrt{\frac{\sum_i k_i}{N}}\text{(stat)} + \frac{\delta}{2}\text{(syst)} 
	\]

Which is a specific case of the more general:

	\[
	\hat\mu(b_0)\pm \sqrt{\Var\left[\hat\mu(b_0)\right]}\pm\frac{\max \hat\mu\left(b_1\leq b \leq b_2\right)-\min\hat\mu\left(b_1\leq b \leq b_2\right)}{2}
	\]

There isn't a convention that is universally accepted, though usually both the systematic and statistic uncertainties get written separately, and each combine with their own counterparts (though systematic uncertainties are also often summed in quadrature like statistic ones). 

\begin{example}
	Let's look at an interesting case: along with my observable $k$, distributed like a Poisson $Pois(k,\mu+b)$, we have an observable $q$ that has a pdf $p(q;b)$. This means that there is \emph{some} observable that gives me information on my background parameter $b$ (from which I was able to determine that $b$ fell within a certain range, for example). 
	
	Well, I can now express my experiment as a combination measurement of two observables
	
	\[
	p(k,q;\mu,b) = Pois(k;\mu+b)\cdot p(q;b)
	\]
	
	from which I can find estimates of the two parameters through techniques like the MLE. 
	
	In this case, the systematic uncertainty gets absorbed in the statistic, though now we have paired $q$ directly to $k$ and we won't be able to use them independently. 
\end{example}

\begin{example}
	Another interesting case is where $b$ has a pdf (or \emph{dob} in the Bayesian case) and we can eliminate it by marginalizing:
	
	\begin{align*}
	p_q(k;\mu) &= \int p(k,b|q;\mu)\dd b = \int \pi(b)\frac{p(k,1|b;\mu)}{p(q)}\dd b\\
	&= \frac{1}{p(q)}\int \pi(b) Pois(k;\mu+b)p(q;b)db
	\end{align*}
	
	Also in this case the systematic uncertainty gets absorbed in the statistic. In the Bayesian case, in which we are free to assign a degree of belief to any parameter, the systematic uncertainty is always treatable like the statistic one. However, it's fundamental that $b$ follows a probability distribution (and thus mus be assigned a \emph{prior}), which is what it means when "the nuisance parameters were integrated" is written.
\end{example}


Of course, we are lucky enough to have Likelihood-like information on $b$ then we can naturally estimate it, using the usual techniques. 

We also take a moment to talk about non-systematic "uncertainties" that might plague our data. Sometimes, we find ourselves in the situation where our data depends on unknown or ambiguous parameters that are not uncertainties, and therefore must not be treated as such. For example:
\begin{itemize}
	\item Histogram binning and range
	\item Number of observables to consider
	\item Choice of subsamples to include
	\item Choice of data cuts and limits
\end{itemize}

Of course, the result of our measurements will depend on these (and so will the pdf of the data), though they are not uncertainties. 

The important thing is to remain transparent at all times on these arbitrary choices, and to show in your treatment of the data that the results are not very sensitive to the different choices one may make (for example, by changing the binning and proving that the results do not vary significantly).

In general, systematic uncertainties are a sensitive and fragile topic, and thus solutions that are less dependent on systematic errors or arbitrary experimental decisions are preferred over ones that are heavily affected by it.